\documentclass[10pt,a4paper]{amsart}
\usepackage[margin=19mm]{geometry}
\usepackage{amsmath,amssymb,mathtools}
\usepackage[scr=rsfs,cal=euler]{mathalfa}
\usepackage{xfrac}
\usepackage[unicode,hidelinks,bookmarks=false]{hyperref}
\newcommand{\ie}{i.e.\ }
\newcommand{\cf}{cf.\ }
\newcommand{\resp}{respectively\ }
\newcommand{\Subsection}{\subsection}
\providecommand{\nobreakdash}{\nobreak\hbox{-}}
\setlength{\emergencystretch}{2em}
\multlinegap0pt
\usepackage{bm}
\usepackage{bbm}
\usepackage{dsfont}
\usepackage{xspace}
\usepackage{cancel}
\usepackage{mathtools}
\usepackage[shortlabels]{enumitem}
\usepackage{amsthm}
\makeatletter
\@ifundefined{theorem}{\newtheorem{theorem}{Theorem}}{}
\@ifundefined{proposition}{\newtheorem{proposition}[theorem]{Proposition}}{}
\makeatother
\makeatletter
\expandafter\ifx\csname Proof\endcsname\relax
\newenvironment{Proof}
{\par\textit{Proof}.}
{\hfill$\square$\par\vspace{0.2cm}}
\fi
\makeatother
\title[On Congruence Theorem for valued division algebras]{On Congruence Theorem for valued division algebras}
\author{Huynh Viet Khanh, Nguyen Duc Anh Khoa}
\date{Source: arXiv:2503.17714v2 (CC BY 4.0)}
\begin{document}
\maketitle

\noindent\emph{Source and licence.} On Congruence Theorem for valued division algebras. Version arXiv:2503.17714v2 (CC BY 4.0), distributed under the Creative Commons Attribution 4.0 International licence, as recorded in the arXiv licence metadata for this exact version and shown on the abstract page https://arxiv.org/abs/2503.17714v2. Full rights notice: licenses/arxiv-2503.17714-CC-BY-4.0.txt.\par\smallskip

\noindent\textit{Editorial scope.} Counted unit: the Proposition whose source label is \texttt{corollary\_Qp} in the article (source0.tex lines 122--128, source label \texttt{corollary\_Qp}), delivered together with the article's own complete original proof of it (source0.tex lines 129--159, one \texttt{proof} environment of 31 lines). No mathematics is written, repaired or re-derived: statement and proof are the article's own text line for line. Required context retained uncounted: Theorem\_p-group (lines 91--93). Every definition, display and label of those lines is retained unchanged. Omitted from this package and not redistributed: 1--90, 94--121, 160--301. Nothing inside a retained excerpt is omitted or altered: every retained body is delivered exactly as the article's lines are written, so no omission receipt of any kind accompanies this package. Citations: the retained excerpts carry no citation command at all, so no bibliography entry of the article is transcribed and no reference is redistributed. Reference adaptation: none was needed and none was made; every reference inside the delivered unit resolves inside it, so no unresolved reference or citation is produced. Printed slips kept literal: no printed word, symbol, hypothesis or number of the counted statement or of its proof was altered. Layout and class: the article's own class is \texttt{amsart} with packages including mathtools, amsfonts, amsthm, amssymb, enumitem, tikz, inputenc, fontenc, geometry, hyperref; the wrapper is the lane's pinned \texttt{amsart} editorial wrapper at 10pt on A4, which restates the theorem-like environments the retained text needs and adds the article's own macro definitions that the retained text needs (none), with extra wrapper packages (bm, bbm, dsfont, xspace, cancel, mathtools, enumitem). The delivered numbering is the unit's own. Assets: no retained span contains a figure environment, a graphics-inclusion command, a PDF-page inclusion, a table or a TikZ picture; no image file is copied into this package, the package declares no asset dependency and no image file is redistributed with it. Licence: version arXiv:2503.17714v2 of this article is distributed under the Creative Commons Attribution 4.0 International licence, as recorded in the arXiv licence metadata for this exact version and shown on the abstract page https://arxiv.org/abs/2503.17714v2. A full rights notice is delivered at licenses/arxiv-2503.17714-CC-BY-4.0.txt. Characters: every retained excerpt is pure ASCII, so the delivered engine, XeLaTeX with its default Latin Modern text font, reports no missing character.
\begin{theorem}\label{Theorem_p-group}
	Let $K$ be a field with Henselian valuation, and let $D$ be a tame $K$-central division algebra. Then, either $(1+M_D) \cap {\bf G} \subseteq D'$ or ${\rm char}(\overline{K})=p>0$ and ${\bf H}$ is a $p$-group.
\end{theorem}

\begin{proposition}\label{corollary_Qp}
	Let $(\mathbb{Q}_p,v_p)$ be the field of $p$-adic numbers with $p$ a prime, and $v_p$ the $p$-adic valuation on $\mathbb{Q}_p$. Let $D$ be a tame ${\mathbb{Q}_p}$-central division algebra with ${\rm ind}(D) = n$. Then, the following assertions hold:
	\begin{enumerate}[font=\normalfont]
		\item If $p>2$ or  $p=2$ and $n$ is odd, then ${\bf H}\cong \{\pm 1\}$.
		\item If $p=2$ and $n$ is even, then ${\bf H}=1$.
	\end{enumerate}
\end{proposition}

\begin{proof}
	We consider three possible cases:
	
	\medskip 
	
	\medskip 
	
	\textit{\textbf{Case 1}: $p>2$. } Then $\tau\left(\mathbb{Q}_p^*\right) = \left\{\varepsilon_1,\varepsilon_2,\ldots,\varepsilon_{p-1}\right\}$, where $\varepsilon_i$ is a $(p-1)$-th root of unity in $\mathbb{Q}_p$ for all $i \in \{1,2,\ldots,p-1\}$. As $\varepsilon_i^{p-1} = 1$, it follows that $\tau\left(\mathbb{Q}_p^*\right)$ contains no element of order $p$. Assume by contradiction that ${\bf H}$ has an element of order $p$, say $aD'$ where $a\in D$. Then, we have $a^p \in D'$ and $a \notin D'$. It follows that ${\rm Nrd}_D(a)^p= {\rm Nrd}_D\left(a^p\right) =1$, which implies that ${\rm Nrd}_D(a) \in \tau(\mathbb{Q}_p^*)$. Since $\tau(\mathbb{Q}_p^*)$ has no elements of order $p$, we get that ${\rm Nrd}_D(a) =1$, and hence $a \in D^{(1)}$. Because $a\in \left(1+M_D\right)$, we get $a \in \left(1+M_D\right)\cap {\bf G} \subseteq D'$, which is a contradiction. According to Theorem \ref{Theorem_p-group}, we conclude that ${\bf H}=1$.
	
	\medskip 
	
	\textit{\textbf{Case 2}: $p=2$ and $n$ is even.} Then $\tau\left(\mathbb{Q}_2^*\right) = \left\{\pm1\right\}$. As ${\rm TK}_1(D)/{\rm SK}_1(D)$ is isomorphic to $ {\rm Nrd}_D\left({\rm TK}_1(D)\right) \le \tau\left(\mathbb{Q}_2^*\right)$, we get that
	$$
	((1+M_D)\cap {\bf G})D^{(1)}/D^{(1)} \cong {\bf H}{\rm SK}_1(D)/{\rm 
		SK}_1(D)  \le 
	\tau\left(\mathbb{Q}_2^*\right).
	$$
	It follows that $((1+M_D)\cap {\bf G})D^{(1)}/D^{(1)} \subseteq \left\{-D^{(1)},D^{(1)}\right\}$. As ${\rm Nrd}_D(-1) = (-1)^n = 1$, we obtain $-D^{(1)} = D^{(1)}$. But we also have $-1 \in\left(1+M_D\right)\cap {\bf G}$, from which it follows that $((1+M_D)\cap {\bf G})D^{(1)}/D^{(1)} = 1$. Hence,
	$$
	\left(1+M_D\right)\cap {\bf G} \subseteq \left(1+M_D\right) \cap D^{(1)} \subseteq D',
	$$
	which implies that ${\bf H}=1$.
	
	\medskip 
	
	\textit{\textbf{Case 3}: $p=2$ and $n$ is odd.} As ${\rm Nrd}_D(-1) = (-1)^n = -1 \ne 1$, it follows that $D^{(1)} \ne -D^{(1)}$. On the other hand, we also have $-1 \in (1+M_D)\cap {\bf G}$. This means that $((1+M_D)\cap {\bf G})D^{(1)}/D^{(1)} = \{D^{(1)},-D^{(1)}\}$. Finally, we have 
	$$
	((1+M_D)\cap {\bf G})D^{(1)}/D^{(1)} \cong {\bf H}{\rm SK}_1(D)/{\rm SK}_1(D) \cong\tau\left(\mathbb{Q}_2^*\right),
	$$
	so ${\bf H} \cong \{\pm 1 \}$.
\end{proof}

\end{document}
